\subsection{Finite-trajectory bounds and fractional innovation for the quartic models}
\label{sec:finite-trajectory-numerics}
\label{sec:fractional-innovation-p4}

We evaluate the estimates of \Cref{sec:c5} on the one-dimensional pure and
regularized \(p=4\), \(k=3\) trajectories, with \(u^*(x)=\cos(2\pi x)\) and
\(\varepsilon=0\) or \(0.1\) (cases \texttt{base\_pure} and \texttt{epsilon\_01}
in Supplement S1). These runs use a larger pool and a smaller accepted-state budget than the
eight cases of \Cref{sec:numerics}. The inner product held fixed during the estimates is
\[
 a_*(v,w)=\int_0^1\big(\varepsilon^2+3|(u^*)'|^2\big)v'w'\dd x.
\]
The calibration range contains the 25 retained states \(N=8,\ldots,32\) and its 24 transitions
\(N\to N+1\), \(N=8,\ldots,31\).  The subsequent finite range contains the 32 transitions
\(N\to N+1\), \(N=32,\ldots,63\), generated by extending the same previously computed iterates to \(N=64\).
The two transition sets are kept separate: all source, comparison, and fractional-innovation
constants are estimated on the calibration range only, while the continuation is used only as a
auxiliary numerical test.
At the starting index, the fixed auxiliary dictionary comprises 463 valid candidates and the first 128
valid reference atoms in the order in which they were generated, drawn with a separate seed. They
need not be linearly independent. Symmetry and the frozen normalization follow the theory, and
Supplement S2 specifies the ordering and the filtering. The reference atoms enter the estimates
only. Source coefficients are constructed from the starting-index residual and are not adjusted to fit
later errors.

The full active-space projections use untruncated QR on a common breakpoint-segmented Gauss
rule of order 20, cross-checked at order 16. The active linear systems have condition numbers
between \(7.7\times10^1\) and \(4.1\times10^3\). The maximum defect in
\cref{eq:c5-quartic} is below \(8.7\times10^{-17}\), and the projection-drop identity
has residual below \(1.5\times10^{-12}\). These are floating-point calculations, not interval
enclosures of the exact integrals or projections.

\paragraph{Source bound and Hessian-transfer conditions.}
\Cref{tab:merged-trajectory-conditions} collects the source data, the quartic comparison
constants, and the initial-value bounds of \Cref{thm:c5-source}, evaluated with
\(\underline\theta_m=\theta_m\), the observed selection ratio. In the table,
\(B=B_{\rm entry}\) and \(q=q_{\rm corr}\). The source energy
bounds at \(N=32\) are \(2.23\times10^{-2}\) and \(2.11\times10^{-2}\), compared with
energy gaps \(1.99\times10^{-6}\) and \(3.01\times10^{-6}\). Thus the source remainder
does not make this bound sharp; the accumulated-innovation contribution remains dominant.
For both models, \(\sigma^2/r_8^2\) is approximately \(6.21\times10^{-5}\), whereas
\(\sigma^2/r_{32}^2\) is \(0.751\) for the pure model and \(0.495\) for the regularized model.

\begin{table}[!htbp]
\centering
\small
\caption{Source and quartic comparison quantities on all retained states \(N=8\)--32. The source is fixed at the starting index; \(q\) and \(\rho\) are calibration-range maxima on \(N=8\)--32. Values are floating-point estimates.}
\label{tab:merged-trajectory-conditions}
\begin{tabularx}{\linewidth}{@{}Xrr@{}}
\toprule
Quantity & Pure \(p=4\) & \(\varepsilon=0.1\) \\
\midrule
Source budget \(B\) & \(4.361\times10^{1}\) & \(4.362\times10^{1}\) \\
Source remainder \(\sigma\) & \(1.726\times10^{-3}\) & \(1.728\times10^{-3}\) \\
\(\mathcal W_{24}\) & \(1.912\times10^{4}\) & \(2.913\times10^{4}\) \\
\(q=\max D_N/r_N\) & \(4.733\times10^{-2}\) & \(4.715\times10^{-2}\) \\
\(\rho=\max |R_N|/(r_N^2+D_N^2)\) & \(2.006\times10^{-2}\) & \(1.997\times10^{-2}\) \\
Computed \(r_{32}^2\) & \(3.969\times10^{-6}\) & \(6.025\times10^{-6}\) \\
Initial-value source bound for \(r_{32}^2\) & \(4.280\times10^{-2}\) & \(4.057\times10^{-2}\) \\
Computed energy gap at \(N=32\) & \(1.985\times10^{-6}\) & \(3.013\times10^{-6}\) \\
Initial-value source energy bound & \(2.231\times10^{-2}\) & \(2.114\times10^{-2}\) \\
\bottomrule
\end{tabularx}
\end{table}


For comparison with the finite-range discussion, the minimum active-space generalized Hessian
eigenvalues relative to the \(H^1\) Gram matrix are \(0.1328\) and \(0.1428\), and the
minimum finite-reference-pool selection ratios are \(0.9876\) and \(0.9911\).
Nevertheless, the sufficient lower bound of \cref{eq:c5-finite-margin} is negative at all 24
transitions in both models. Positive active-space spectra and finite-reference-set ratios therefore do
not establish the continuous relative-score condition. The selected directions remain valid for the
finite numerical experiment, but the transfer proposition cannot be invoked from these quantities alone.
The direct projection estimates below use the selected directions themselves and do not require that
sufficient lower bound to be positive.
The decomposition in \Cref{tab:transfer-margin-example} shows one representative transition.  The Hessian-transfer lower bound is negative even though the corresponding transition is retained; the score-quadrature and active-optimization terms are negligible at this step relative to the nonlinear-to-frozen defect.
\begin{table}[!htbp]
\centering
\scriptsize
\caption{One-step error-budget decomposition at the first calibration transition. The negative margin is a diagnostic: the nonlinear-to-frozen defect dominates the available frozen decrease even though the native step is accepted.}
\label{tab:transfer-margin-example}
\begin{tabularx}{\linewidth}{@{}Xrr@{}}
\toprule
Quantity & Pure \(p=4\), \(8\to9\) & Regularized \(\varepsilon=0.1\), \(8\to9\) \\
\midrule
Frozen decrease & \(1.8403\times10^{-2}\) & \(1.8175\times10^{-2}\) \\
Nonlinear-to-frozen defect & \(6.8347\times10^{-3}\) & \(6.8108\times10^{-3}\) \\
Score quadrature defect & \(9.73\times10^{-13}\) & \(9.95\times10^{-13}\) \\
Active optimization defect & \(4.36\times10^{-7}\) & \(4.48\times10^{-7}\) \\
Total score defect & \(6.8351\times10^{-3}\) & \(6.8112\times10^{-3}\) \\
Selection margin & \(-6.3599\times10^{-2}\) & \(-6.2742\times10^{-2}\) \\
\bottomrule
\end{tabularx}
\end{table}


\paragraph{Fractional-innovation calculation.}
For the benchmark \(\alpha=6\), fixed before evaluating the normalized drops, set
\begin{equation}
 x_N=r_N^2,\qquad
 \gamma_N=\sqrt{x_N-x_{N+1}},\qquad
 c_N=\frac{x_N-x_{N+1}}{x_N^{7/6}},\qquad N=8,\ldots,31.
 \label{eq:numerical-fractional-ratio}
\end{equation}
The minimum over all transitions is \(1.6904\times10^{-2}\) for the pure model and
\(1.0912\times10^{-2}\) for \(\varepsilon=0.1\); both occur at \(25\to26\).
Changing the quadrature order changes these stepwise ratios by less than
\(5.6\times10^{-12}\) relatively. Together with the computed \(q\) and \(\rho\), these
values permit a numerical evaluation of \cref{eq:c5-power,eq:c5-power-energy}.
Writing \(j=N-8\), the resulting finite-range comparison powers are
\[
 \widehat\delta_{8+j}=O(j^{-6}),\qquad
 d_{V,\bullet}(\widehat G_{8+j},u^*)=O(j^{-3}),\qquad
 \|\widehat G_{8+j}-u^*\|_{W^{1,4}}=O(j^{-3/2}),\quad 1\leq j\leq24.
\]
The explicit bound in \cref{eq:c5-power} also includes the starting state \(j=0\).
These powers describe a conditional envelope on the tested finite range, not a limit as \(N\to\infty\).

\begin{table}[!htbp]
\centering
\small
\caption{Fractional-innovation quantities for \(\alpha=6\), computed from every transition \(N\to N+1\), \(N=8,\ldots,31\). The fitted exponent describes the innovation score and is not used to choose the benchmark exponent or its bound.}
\label{tab:fractional-p4}
\begin{tabularx}{\linewidth}{@{}Xrr@{}}
\toprule
Quantity & Pure \(p=4\) & \(\varepsilon=0.1\) \\
\midrule
\(c_{\min}\) over all 24 transitions & \(1.690\times10^{-2}\) & \(1.091\times10^{-2}\) \\
Transition attaining \(c_{\min}\) & \(25\to26\) & \(25\to26\) \\
Fitted slope \(\nu\) of \(\log\gamma_N\) versus \(\log x_N\) & 0.5866 & 0.5898 \\
\(1/(2\nu-1)\) from the fit & 5.77 & 5.57 \\
Fractional bound for \(r_{32}^2\), \(\alpha=6\) & \(3.774\times10^{-2}\) & \(4.110\times10^{-2}\) \\
Fractional bound / computed \(r_{32}^2\) & \(9.510\times10^{3}\) & \(6.821\times10^{3}\) \\
Fractional energy bound at \(N=32\) & \(1.967\times10^{-2}\) & \(2.142\times10^{-2}\) \\
\bottomrule
\end{tabularx}
\end{table}


A least-squares fit of \(\log\gamma_N\) against \(\log x_N\), using all 24 transitions,
gives slopes \(\nu=0.5866\) and \(0.5898\). The relation
\(\alpha=1/(2\nu-1)\) associates these fitted innovation powers with \(5.77\) and
\(5.57\), respectively. This fit concerns \(\gamma_N=\theta_NS_N/\ell_N\), not \(S_N\)
alone; it does not independently establish either factor in \cref{eq:c5-power-factors}.

The worst-step constants still yield conservative terminal projection bounds:
\(3.77\times10^{-2}\) and \(4.11\times10^{-2}\), approximately
\(9.5\times10^3\) and \(6.8\times10^3\) times the computed projection errors.
The fractional bound is slightly smaller than the source bound for the pure model and slightly
larger for the regularized model. Positive drops on a finite sequence also give a positive
minimum normalized ratio for any fixed exponent. The computation establishes the stated numerical
envelope and its constants, but does not uniquely identify sixth-order decay or verify an
asymptotic mechanism. Such a conclusion would require bounds uniform over increasing index ranges.

\paragraph{Continuation beyond the calibration range.}
To test whether the calibrated constants remain informative immediately beyond the range on which
they were fixed, we continued the two trajectories from \(N=32\) to \(N=64\), keeping the same
candidate pools, previously computed iterates, and source coefficients and drawing no further random atoms. The
comparison exponent and all constants were kept at their \(N=8\)--\(32\) values. Most transitions in
this subsequent finite range satisfy the normalized-innovation sufficient condition with those constants.
At quadrature order \(20\), it holds for 27 of the 32 continuation transitions in the pure model
\(27/32=84.4\%\) and 28 of 32 in the regularized model
\(28/32=87.5\%\); only five and four transitions, respectively, fall
below the prescribed threshold. These exceptions concern the sufficient condition, not the validity
of the theoretical implication. In both models, the first such transition is \(N=33\to34\), and the
minimum ratio \(c_N/c_0\) over the continuation is \(0.156\) and \(0.187\), respectively; the
terminal squared projection errors are \(4.770\times10^{-7}\) and \(4.626\times10^{-7}\). The
order-16 cross-check changes these reported values only in the displayed last digits. The
continuation therefore shows that the calibrated condition remains satisfied for
most transitions, together with a substantial error decrease and agreement with the algebraic
identities on the subsequent states. The few exceptions prevent applying the same fixed-constant
theorem over the entire continuation range. The stated finite-range envelope therefore retains
its calibration range $N=8$--32, while the continuation provides supporting stepwise numerical
evidence without establishing an asymptotic rate.

\begin{table}[!htbp]
\centering\small
\caption{Continuation from $N=32$ to $N=64$ (32 transitions) with constants fixed on $N=8$--$32$. Pass counts refer to the normalized-innovation sufficient condition; the first failed transition starts at the reported $N$.}
\label{tab:continuation-window}
\begin{tabular}{@{}lrrrrr@{}}
\toprule
Model & New steps & Pass & First failed $N$ & $\min c_N/c_0$ & $r_{64}^2$ \\
\midrule
Pure $p=4$ & 32 & 27 & 33 & 0.156 & $4.770\times10^{-7}$ \\
$\varepsilon=0.1$ & 32 & 28 & 33 & 0.187 & $4.626\times10^{-7}$ \\
\bottomrule
\end{tabular}
\end{table}

\FloatBarrier
